% from ega0-19b-fr.tex
\oldpage[0\textsubscript{IV-1}]{113}

dans les deux cas, on a $\dim(B)=\dim(A)=m$~; dans le premier cas, cela résulte
de (17.1.4, (iii))~; dans le second, on voit directement que $p$ et les $T_i$
($1\leq i\leq m-1$) forment une suite $B$-régulière engendrant $\mathfrak n$,
ou on peut aussi utiliser le fait que ces éléments engendrent $\mathfrak n$ et
que l'on a $\dim(B)\geq\dim(A)$ par (16.3.10). Comme $A$ et $B$ sont intègres,
on tire finalement de (16.3.10) que $w$ est injectif, ce qui achève la
démonstration.
\end{proof}

\begin{corollary}[19.8.9]
\label{0.19.8.9-fr}
Soit $A$ un anneau local noethérien intègre complet contenant un corps $k_0$~;
soit $k$ le corps résiduel de $A$, et supposons que $k$ soit fini sur $k_0$.
Alors, dans la conclusion de (19.8.8, (ii)), on peut remplacer $1^{\mathrm o}$
et $2^{\mathrm o}$ par la condition que $B$ est de la forme
$k_0[[T_1,\ldots,T_n]]$, l'injection canonique $B\to A$ étant un
$k_0$-homomorphisme local (pour la structure usuelle de $k_0$-algèbre de $B$).
\end{corollary}

\begin{proof}
En effet, en reprenant la démonstration de (19.8.8, (ii)), on définit cette fois
$w:k_0[[T_1,\ldots,T_n]]\to A$ comme coïncidant dans $k_0$ avec l'identité et
appliquant $T_j$ sur $y_j$ pour $1\leq j\leq n$. L'hypothèse que $k$ est de
degré fini sur $k_0$ entraîne encore que $A$ est un $B$-module quasi-fini, donc
de type fini par ($\mathbf 0_{\mathrm I}$, 7.4.1), et on conclut comme dans
(19.8.8).
\end{proof}

\begin{corollary}[19.8.10]
\label{0.19.8.10-fr}
Soit $A$ un anneau local artinien dont l'idéal maximal $\mathfrak m$ est de
carré nul~; il existe alors un anneau local noethérien régulier $B$, d'idéal
maximal $\mathfrak n$, tel que $A$ soit isomorphe à $B/\mathfrak n^2$.
\end{corollary}

\begin{proof}
Soient $K$ le corps résiduel $A/\mathfrak m$ de $A$, $n$ le rang de
$\mathfrak m/\mathfrak m^2=\mathfrak m$ sur $K$. Si $A$ contient un corps, il
résulte de (19.6.3) que $A$ est isomorphe à $B/\mathfrak b$, où
$B=K[[T_1,\ldots,T_n]]$ et $\mathfrak b$ est contenu dans le carré de l'idéal
maximal $\mathfrak n$ de $B$~; mais comme
$\operatorname{long}(B/\mathfrak n^2)=n+1=\operatorname{long}(A)$, on a
nécessairement $\mathfrak b=\mathfrak n^2$.

Supposons ensuite que $A$ ne contienne pas de corps~; cela entraîne que $K$ est
de caractéristique $p>0$ et que $p\cdot1\neq0$ dans $A$ (19.6.3)~; donc
$p\cdot1$ est un élément de $\mathfrak m$, et il y a par suite $n-1$ autres
éléments $x_i$ ($2\leq i\leq n$) de $\mathfrak m$ formant avec $p\cdot1$ une
base de $\mathfrak m$ sur $K$. Soit $W$ un anneau de Cohen dont le corps
résiduel est isomorphe à $K$~; $W$ est un anneau de valuation discrète dont $p$
engendre l'idéal maximal~; on a vu dans la démonstration de (19.8.8) qu'il y a
un homomorphisme $u:W\to A$ appliquant $p$ sur lui-même et qui par passage aux
quotients donne l'identité sur $K$. On prend $B=W[[T_2,\ldots,T_n]]$ et l'on
considère l'homomorphisme local $w:B\to A$ qui coïncide avec $u$ dans $W$ et
est tel que $w(T_i)=x_i$ pour $2\leq i\leq n$. Il est clair que $w$ est
surjectif et que son noyau $\mathfrak b$ est contenu dans le carré de l'idéal
maximal $\mathfrak n=pB+BT_2+\cdots+BT_n$ de $B$~; comme
$\operatorname{long}(B/\mathfrak n^2)=n+1=\operatorname{long}(A)$, on a encore
$\mathfrak b=\mathfrak n^2$.
\end{proof}

\begin{proposition}[19.8.11]
\label{0.19.8.11-fr}
Soient $A$ un anneau local artinien, $\mathfrak m$ son idéal maximal, $k$ son
corps résiduel. Pour que $A$ soit isomorphe à un anneau quotient d'un anneau de
Cohen, il faut et il suffit que $\mathfrak m$ soit engendré par $p\cdot1$, où
$p$ est la caractéristique de $k$.
\end{proposition}

\begin{proof}
La condition est évidemment nécessaire (19.8.5). Pour voir qu'elle est
suffisante, on observe, comme au début de la démonstration de (19.8.8), qu'il
existe un anneau de Cohen $W$ dont le corps résiduel est isomorphe à $k$ et un
homomorphisme local $u:W\to A$. En outre, si l'on considère l'homomorphisme
composé $\mathbf Z_{p\mathbf Z}\to W\xrightarrow{u}A$ (qui est nécessairement
l'unique homomorphisme de $\mathbf Z_{p\mathbf Z}$ dans $A$), on voit que
l'image

